\section{Assembly}\label{sec:assembly}

We use the two seeds through the following equal-norm formulation.

\begin{proposition}\label{prop:equalrow}
There is an absolute constant $C_0$ such that every $X\in\R^{n\times d}$ with
$1\le d$, $2d\le n$ and $\norm{x_i}^2=a$ for all $i$ admits an ENP frame $W$
with
\[
 \fro{X-W}^2\le C_0\,d\,\opn{X^TX-I}.
\]
\end{proposition}

\begin{proof}
Let $\eta=\opn{X^TX-I}$; if $\eta=0$ take $W=X$. Let $A,\eps_0,C_m$ be the
constants of \cref{thm:moderate} and $B,c_*,C_*$ those of \cref{thm:manyrow},
and fix an integer $d_0\ge2$ with $A\log(2Bd^2)\le d$ for all $d\ge d_0$.

\emph{Case $d<d_0$.} \Cref{lem:HM} gives $\fro{X-W}^2\le\eta d(d-1)\le d_0\eta d$.

\emph{Case $d\ge d_0$, $n\ge Bd^2$.} If $\eta\le c_*$ use \cref{thm:manyrow};
otherwise \cref{cor:exist} gives cost $4d\le(4/c_*)\eta d$.

\emph{Case $d\ge d_0$, $n<Bd^2$.} Then $A\log(2n)\le d$. If $\eta\le\eps_0/4$, let
$U=X(X^TX)^{-1/2}$. The singular values of $X$ lie in
$[\sqrt{1-\eta},\sqrt{1+\eta}]$, so $\fro{X-U}^2\le d\eta^2$, and
$\norm{u_i}^2=x_i^T(X^TX)^{-1}x_i\in[a/(1+\eta),a/(1-\eta)]$, so
$|\norm{u_i}^2-a|\le2\eta a$. \Cref{thm:moderate} with $\eps=2\eta$ gives an ENP
$W$ with $\fro{U-W}^2\le2C_m\eta d$, and
$\fro{X-W}^2\le2d\eta^2+4C_m\eta d$. If $\eta>\eps_0/4$ use the trivial cost
$4d\le(16/\eps_0)\eta d$. Altogether
$C_0=\max\{d_0,C_*,4/c_*,2+4C_m,16/\eps_0\}$ works.
\end{proof}

\begin{proof}[Proof of \cref{thm:projection}]
Replacing $(P,Q)$ by $(I-P,I-Q)$ preserves $\beta$, ranks of complements and
distances, so we may assume $d\le n/2$; the case $d=0$ is trivial. Write
$P=UU^T$ with $U^TU=I$, $p_i=\norm{u_i}^2$ and $\delta=\beta/a$. If
$\delta\ge\frac12$, any ENP $W$ (\cref{cor:exist}) gives
$\fro{P-WW^T}^2\le2d\le4n\beta$. Otherwise let $x_i=\sqrt a\,u_i/\norm{u_i}$. Then
$\fro{X-U}^2=\sum_i(\sqrt{p_i}-\sqrt a)^2\le\sum_i(p_i-a)^2/a\le\delta^2d$ and
$X^TX=U^T\Diag(a/p)U$, so $\opn{X^TX-I}\le\max_i|a/p_i-1|\le\delta/(1-\delta)\le
2\delta$. \Cref{prop:equalrow} gives an ENP $W$ with
$\fro{U-W}^2\le2\delta^2d+4C_0\delta d\le(1+4C_0)n\beta$, using $\delta d=n\beta$.
Take $Q=WW^T$; by \cref{lem:align}, $\fro{P-Q}^2\le2(1+4C_0)n\beta$. So
\cref{thm:projection} holds with $C_P=\max\{4,2+8C_0\}$.
\end{proof}

\begin{proof}[Proof of \cref{thm:main}]
If $\eps\ge\frac12$, an ENP $W$ gives $\fro{X-W}^2\le2\fro X^2+2d\le6d\le12\eps
d$, since $\fro X^2=\tr X^TX\le(1+\eps)d$. Let $\eps<\frac12$ and
$U=X(X^TX)^{-1/2}$, $P=UU^T$. As above $\fro{X-U}^2\le\eps^2d$ and
$P_{ii}=x_i^T(X^TX)^{-1}x_i\in[\frac{1-\eps}{1+\eps}a,\frac{1+\eps}{1-\eps}a]$, so
$\beta=\max_i|P_{ii}-a|\le4\eps a$. \Cref{thm:projection} gives $Q$ with
$\fro{P-Q}^2\le4C_P\eps d$, and by \cref{lem:align} there is an isometry $W$
spanning $Q$---hence ENP---with $\fro{U-W}^2\le\fro{P-Q}^2$. Then
$\fro{X-W}^2\le2\eps^2d+8C_P\eps d$, so \cref{thm:main} holds with
$C=\max\{12,1+8C_P\}$.
\end{proof}
